% Folder 34, batch 10 (archivists' pages 181-200).
% Pass: Opus 5.5 (claude-opus-5-5), 2026-10-03 — first pass, unchecked against the pages by a human.
%
% Page 181 is blank and absent. Pages 182-187 carry his pagination II.31-II.36.
\documentclass[11pt,a4paper]{article}
\input{../preamble/grothendieck}

\folder{34}
\batch{10}
\pages{181}{200}
\foldertitle{SGA 7 : notes manuscrites (s.d.).}
\dating{[vers 1967-1973]}
\watermark{Édition de démonstration}

\begin{document}

\page{182}
\note{p. II.31 de l'auteur ; la section 6 du chapitre II commence ici, les sections précédentes sont hors de ce lot}

\section{\struck{Appendice.} 6. Groupes \add{et algèbres} de Lie associés aux représentations $\ell$-adiques de groupes pro-finis}

\marginal{au-dessus du titre, rattaché à « Groupes » par un trait : analytiques \struck{ou} algébriques etc.}

On a fixé un nombre premier $\ell$.

6.1. Rappelons que si un groupe topologique admet une structure de groupe
analytique sur $\mathbb{Q}_\ell$, celle-ci est \uncertain{essentiellement} unique
(\ill{} si le groupe \ill{} \uncertain{discret}, \ill{} un \ill{} déterminé) ;
\struck{\ill{}} par exemple, \add{tout} homom. continu de groupes analytiques
sur $\mathbb{Q}_\ell$ est automatiquement analytique. En fait, un s.-groupe fermé
d'un groupe analytique (sur $\mathbb{Q}_\ell$) est encore analytique
\uncertain{d'une façon unique}. \ill{}

6.2. Soit $\Pi$ un \ill{} \uncertain{pro-fini}\ill{}. Si
\[
  G = \Pi/V \quad\text{et}\quad G' = \Pi/V'
\]
sont analytiques sur $\mathbb{Q}_\ell$, alors
\[
  G'' = G/V\cap V'
\]
aussi, car il \struck{\ill{}} est isomorphe à un s.-groupe fermé de $G\times G'$\note{on attend $\Pi/V\cap V'$ ; la lettre initiale est écrite comme le $G$ de la ligne précédente}.
Donc les \struck{\ill{}} quotients \emph{analytiques} de $G$\note{sic, pour $\Pi$} forment un
système projectif de groupes analytiques
\[
  \Pi^{\mathrm{an}} = (G_\alpha)
\]
\note{$\Pi^{\mathrm{an}}$ est récrit au-dessus d'une notation biffée, illisible, de la forme « $\ast(\Pi)$ »}
qui \uncertain{est} \ill{} \uncertain{système} projectif de \uncertain{tous} les quotients
\ill{}. Notons que si \ill{} considère $\pi$ \uncertain{comme} \struck{\ill{}}
\add{pro-fini, considéré comme} pro-objet de groupes \struck{analytiques}
\ill{}, \ill{} un \ill{} évidemment
\[
  \Pi^{\mathrm{an}} \longrightarrow \Pi .
\]
\note{devant la flèche, un symbole biffé illisible}

6.3. Les représentations $\ell$-adiques de $\Pi$ (sur $\mathbb{Q}_\ell$)
\ill{} \uncertain{choisissant} \ill{} \add{continues (= analytiques)} comme
représentations \struck{\ill{}} de $\mathcal{S}_\ell(\Pi)$,
\note{la lettre initiale de « $\mathcal{S}_\ell(\Pi)$ » est douteuse ; c'est vraisemblablement l'ancien nom de $\Pi^{\mathrm{an}}$}
i.e. la catégorie \add{\uncertain{des}} des représentations $\ell$-adiques de $\Pi$ est
\struck{\ill{}} \add{$\varinjlim$} des catégories analogues

\page{183}
\note{p. II.32 de l'auteur}
\uncertain{formées} avec les $G_\alpha$, qui sont \uncertain{des} conditions \uncertain{à
nature} analytique.

6.4. Le syst. projectif des $G_\alpha$ définit un système projectif d'algèbres
de Lie $\mathfrak{g}_\alpha$ sur $\mathbb{Q}_\ell$,
\[
  \text{\struck{\ill{}}}\quad \mathfrak{g} = \mathrm{Lie}^{\mathrm{an}}(\Pi) = (\mathfrak{g}_\alpha)
\]
et \uncertain{on a} un foncteur \struck{$\mathrm{Mod}\,G$}
\[
  \mathrm{Mod}_{\mathbb{Q}_\ell}(\Pi) \longrightarrow \mathrm{Mod}\,\mathfrak{g}
\]
\note{sous la source : « catégorie des modules $\ell$-adiques sur $\Pi$ (sur $\mathbb{Q}_\ell$) » ;
sous le but : « $\varinjlim$ catégorie des modules \struck{\ill{}} $\mathfrak{g}_\alpha$ de dim. finie sur $\mathbb{Q}_\ell$ »}

Ce foncteur est fidèle, mais en général pas \emph{pleinement fidèle} : on
\ill{} des \uncertain{données} \ill{} (\ill{}, on \uncertain{perd} \ill{}). Mais soit
$(U_i)$ les \struck{\ill{}} s.-groupes \uncertain{ouverts} de $\Pi$. $U_i \to \Pi$
définit un \emph{isom.}
\[
  \mathfrak{g}(U_i) \simeq \mathfrak{g}(\Pi),
\]
d'où un foncteur
\[
  \mathrm{Mod}_{\mathbb{Q}_\ell}(U_i) \longrightarrow \mathrm{Mod}\,\mathfrak{g}
\]
et par passage à la limite :
\[
  \boxed{\ \varinjlim_i \mathrm{Mod}_{\mathbb{Q}_\ell}(U_i) \longrightarrow \mathrm{Mod}\,\mathfrak{g}\ }
\]
Cette flèche est une équivalence de \emph{catégories}.

\page{184}
\note{p. II.33 de l'auteur}
Cela se \uncertain{ramène} facilement \struck{au cas où \ill{} $\pi$ \ill{}}
\add{\ill{}} \struck{analytique \ill{}} au

\emph{Lemme}. Soit $G$ un groupe analytique, $\mathfrak{g}$ son alg. de Lie, qui
est aussi l'alg. de Lie de ses sous-groupes ouverts $U_i$. Alors le foncteur
\[
  \varinjlim_i \mathrm{Mod}_{\mathbb{Q}_\ell} U_i \longrightarrow \mathrm{Mod}_{\mathbb{Q}_\ell}\,\mathfrak{g}
\]
est une équivalence de catégories.

6.5. On fera attention \struck{\ill{}} que même si $\pi = \hat{\mathbb{Z}}$,
\add{\uncertain{que} \ill{}} \struck{\ill{}} on ne récupère pas $\pi$ (\uncertain{ou}
\ill{} \uncertain{plutôt} $\mathfrak{g}$) par la catégorie
$\mathrm{Mod}_{\mathbb{Q}_\ell}(\pi)$, même munie de sa structure de produit
tensoriel : [Car \ill{} \uncertain{ne} \ill{} que \ill{} \uncertain{algébriques}
\struck{$G$} \ill{} les $\mathfrak{g}_\alpha$ \ill{} \ill{} \ill{} …].

\marginal{dans un cadre, à gauche de 6.6 : Attention, il est essentiel ici de travailler sur $\mathbb{Q}_\ell$ ! \ill{} \ill{} \uncertain{remarque} !}

\struck{6.6.} Je vais construire un pro-groupe \emph{\add{pro-}algébrique} \add{affine}
$\pi^{\mathrm{alg}}$ sur $\mathbb{Q}_\ell$\note{« $\mathbb{Q}_\ell$ » est cerclé},
\struck{$\pi$} \uncertain{canonique}, \uncertain{tel} que $\mathrm{Mod}_{\mathbb{Q}_\ell}(\pi)$
\uncertain{soit} équivalente \ill{} structure tensorielle. Pour \uncertain{cela},
\ill{} \uncertain{considérer}, pour tout $E \in \mathrm{Ob}\,\mathrm{Mod}_{\mathbb{Q}_\ell}(\pi)$,
le groupe \uncertain{algébrique} engendré dans $\underline{\mathrm{Aut}}_{\mathbb{Q}_\ell}(E)$
par l'image de $\pi$, soit $\widetilde{G}_E$. S'il existe un homom.
\emph{surjectif} $E \to F$, il y a un homom. naturel $\widetilde{G}_E \to \widetilde{G}_F$
\add{qui est un épimorphisme (\uncertain{surjectif})}, et de cette façon les $(\widetilde{G}_E)$

\page{185}
\note{p. II.34 de l'auteur}
forment un système projectif de groupes algébriques, noté
$\widetilde{G} = \pi^{\mathrm{alg}} = (\widetilde{G}_E)$. On a enfin un foncteur évident
\[
  \mathrm{Mod}(\pi^{\mathrm{alg}}) \xleftarrow{\ \approx\ } \mathrm{Mod}_{\mathbb{Q}_\ell}(\pi)
\]
\note{sous $\mathrm{Mod}(\pi^{\mathrm{alg}})$ : « modules de représentations \emph{algébriques} de $\pi^{\mathrm{alg}}$ » ; la flèche porte une pointe à chaque bout}
qui est une équivalence de catégories. On a d'ailleurs
\[
  \pi^{\mathrm{alg}}/(\pi^{\mathrm{alg}})^{0} \simeq \pi
\]
ce qui signifie que, si $\pi^{\mathrm{alg}} = (\widetilde{G}_\alpha)$
\add{\uncertain{indexé}}, \ill{}
\[
  \pi \xleftarrow{\ \approx\ } (\widetilde{G}_\alpha/\widetilde{G}_\alpha^{0})
\]
(isom. canonique), où bien \uncertain{sûr} on identifie\marginal{cerclé : $G$ algébrique} des groupes finis à des
groupes alg. sur $\mathbb{Q}_\ell$.

6.6.\note{numéro surchargé} Au système projectif $\pi^{\mathrm{alg}} = (\widetilde{G}_\alpha)$
\uncertain{correspond} une pro-algèbre de Lie
\[
  \widetilde{\mathfrak{g}} = \mathrm{Lie}^{\mathrm{alg}}(\pi) = \mathrm{Lie}(\pi^{\mathrm{alg}}) = (\mathrm{Lie}(\widetilde{G}_\alpha)),
\]
et on a un foncteur évident
\[
  \mathrm{Mod}_{\mathbb{Q}_\ell}(\pi) \simeq \mathrm{Mod}(\pi^{\mathrm{alg}}) \longrightarrow \mathrm{Mod}(\widetilde{\mathfrak{g}}),
\]
\note{sous la flèche : « $\varinjlim$ des cat. de modules sur $\widetilde{\mathfrak{g}}_\alpha$ de dim. finie sur $\mathbb{Q}_\ell$ » ;
un argument biffé suit $\widetilde{\mathfrak{g}}$, et le second $\pi^{\mathrm{alg}}$ est cerclé}
\struck{et \ill{}} et \uncertain{donc}
\[
  \varinjlim_{U_i} \mathrm{Mod}_{\mathbb{Q}_\ell}(U_i) \xrightarrow{\ \approx\ } \mathrm{Mod}(\widetilde{G}^{0}) \longrightarrow \mathrm{Mod}(\widetilde{\mathfrak{g}})
\]
\note{sous la limite : « prenant les sous-groupes ouverts des $U_i$ » (indice $U_i$ surchargé) ; sous $\mathrm{Mod}(\widetilde{G}^{0})$ : « syst. projectif des composantes neutres $\widetilde{G}_\alpha^{0}$ des $\widetilde{G}_\alpha$ »}

\page{186}
\note{p. II.35 de l'auteur}
La première flèche est une \emph{équivalence} \add{\uncertain{entre} représentations
$\ell$-adiques} de catégories : le fait de \uncertain{remplacer} $\pi$ par un
s.-groupe ouvert \ill{} \uncertain{revient} à \uncertain{remplacer}
$\widetilde{G} = \pi^{\mathrm{alg}}$ par la \emph{composante neutre}
$\widetilde{G}^{0} = (\widetilde{G}_\alpha^{0})$. (N.B. Je \ill{} \ill{} question de
\uncertain{passer} de la composante \ill{} dans les 2 contextes \ill{} analytique).
Par contre, le foncteur
\[
  \mathrm{Mod}(\widetilde{G}^{0}) \longrightarrow \mathrm{Mod}(\widetilde{\mathfrak{g}})
\]
est seulement \emph{pleinement fidèle}, mais \struck{\ill{}} pas
\struck{\ill{}} \uncertain{essentiellement} surjectif a priori.

6.6. \struck{Nous} \uncertain{Pour} un groupe algébrique $\Gamma$ sur
$\mathbb{Q}_\ell$\add{\uncertain{affine}}, \ill{} aussi $\Gamma$ \ill{}
$\Gamma(\mathbb{Q}_\ell)$. \ill{} \ill{} \uncertain{éléments} \ill{} \ill{}
$\Gamma/\Gamma^{0}$ \uncertain{sont} \ill{} sur $\mathbb{Q}_\ell$, donc
$\Gamma(\mathbb{Q}_\ell)$ \ill{} \ill{} \ill{} \ill{} \ill{} par les groupes
algébriques \uncertain{envisagés} dans \ill{} sections précédentes.

\uncertain{Ceci} dit, nous \uncertain{avons} un homom. canonique
\[
  (6.6.1)\qquad
  \begin{array}{ccc}
    \pi^{\mathrm{an}} & \longrightarrow & \pi^{\mathrm{alg}} \\
    \| & & \| \\
    G & & \widetilde{G}
  \end{array}
\]
qui s'explicite ainsi :
\[
  (\ast)\qquad \pi^{\mathrm{an}} \xrightarrow{\ \approx\ } \text{“}\varprojlim_{E \in \mathrm{Ob}\,\mathrm{Mod}_{\mathbb{Q}_\ell}(\pi)}\text{”}\, G_E
\]
\note{la flèche est doublée d'une flèche de retour}
où $G_E$ est l'image de $\pi$ dans $\mathrm{Aut}_{\mathbb{Q}_\ell}(E)$, et

\page{187}
\note{p. II.36 de l'auteur}
\[
  \pi^{\mathrm{alg}} = \text{\struck{\ill{}}}\ \text{“}\varprojlim\text{”}\, \widetilde{G}_E
\]
où $\widetilde{G}_E$ est l'\emph{enveloppe algébrique} du groupe
\uncertain{analytique} $G_E$.

\emph{NB} En fait, le \uncertain{foncteur} \ill{} est une \emph{bijection}, \ill{}
\uncertain{les} groupes analytiques \uncertain{envisagés} sur $\mathbb{Q}_\ell$
\ill{} une représentation \uncertain{linéaire} fidèle …\marginal{réf ?}

Donc $\pi^{\mathrm{an}} \to \pi^{\mathrm{alg}}$ est un \emph{\add{dense}}
\struck{\ill{}} \add{\uncertain{homom.}} \add{de pro-groupes} \ill{} $G_E$ sont \add{denses}
\struck{\ill{}} \ill{} algébriques.\marginal{?}
\struck{\ill{} jamais \ill{} ; \ill{} $\pi \leftarrow \pi^{\mathrm{alg}}$ \ill{} \ill{}
\ill{} \ill{} \ill{} …}\note{passage de trois lignes biffé et barré d'un trait sinueux ; on y lit « $\pi \simeq \pi^{\mathrm{an}}$ » et « p. ex. »}

\uncertain{Lie} (6.6.1) \uncertain{induit} un homom.
\[
  (6.6.2)\qquad \mathfrak{g} = \mathrm{Lie}^{\mathrm{an}}(\pi) \longrightarrow \widetilde{\mathfrak{g}} = \mathrm{Lie}^{\mathrm{alg}}(\pi) .
\]
C'est un isomorphisme si et seulement si les alg. de Lie
$\mathfrak{g}_E = \mathrm{Lie}(G_E)$ sont \emph{algébriques}, \uncertain{i.e.} si
\ill{} \struck{\ill{}} \add{\ill{}} les $G_E$ sont \uncertain{ouverts} dans les
$\widetilde{G}_E$. \struck{Dans ce cas, le \ill{}} \ill{} \uncertain{foncteur}
\[
  (6.6.3)\qquad \varinjlim_{\substack{U_i \text{ ouverts}\\ \text{dans } \pi}} \mathrm{Mod}_{\mathbb{Q}_\ell}(U_i) \xrightarrow{\ \approx\ } \mathrm{Mod}(\widetilde{\mathfrak{g}})
\]
\ill{} \uncertain{donc} une équivalence de catégories ….

\marginal{laissons tomber remarques 6.7. et 6.8. dans un \ill{}}
\emph{Remarques} 6.7. Si \ill{} \ill{} la \uncertain{première} \ill{} \ill{}
$\mathrm{Mod}(\pi^{\mathrm{alg}})$, l'équivalence \ill{} (6.6.3) \ill{}

\page{188}
\note{p. II.37 de l'auteur}
peut surprendre, car \uncertain{en} \uncertain{supposant} que les $\widetilde{G}_E$
soient semi-simples disons) \struck{\ill{} \ill{} qu'ils soient} ils ne sont
pas nécessairement simplement connexes (au sens de Chevalley) i.e. ils
pourraient avoir (\uncertain{beaucoup} …) moins de représentations que leurs
algèbres de Lie. L'explication provient du fait que dans le système projectif
des $\widetilde{G}_{E'}$ majorant $\widetilde{G}_E$, \uncertain{interviennent} aussi les
« revêtements simplement connexes » de $\widetilde{G}_E$, au sens de Chevalley.

\struck{Pour} \uncertain{S'en} \uncertain{convaincre} directement, \struck{\ill{} \ill{}}
\struck{soit $\pi_E$ \ill{}} \add{soit} $G'$ le rev. simplement connexe de
$\widetilde{G}_E^{0}$, et soit $\Gamma$ l'image de \struck{$G' \to \widetilde{G}_E$}
\[
  G'(\mathbb{Q}_\ell) \longrightarrow \widetilde{G}_E^{0}(\mathbb{Q}_\ell).
\]
C'est un s.-groupe distingué de $G'$, \uncertain{donc} il existe un s.-groupe
\uncertain{ouvert} $U$ de $G'(\mathbb{Q}_\ell)$ tel que
$U \to \widetilde{G}_E^{0}(\mathbb{Q}_\ell)$ \ill{} \struck{\ill{}} \uncertain{restriction}
injective \uncertain{\ill{}}. \ill{} \ill{}). Si $E'$ est la \uncertain{source} d'une
représentation \add{\uncertain{(alg.)}} \uncertain{linéaire} fidèle de $G'$, il induit
une représentation \add{analytique} de $U_\ell$, donc une représentation \ill{} d'un
$\widetilde{G}_{E'}^{0}$ ….
\marginal{entre crochets : laissons tomber}

6.8. Si $G$ est un groupe algébrique \struck{affine}
\add{\uncertain{quelconque} (\uncertain{quotient} \ill{})} défini sur $\mathbb{Q}_\ell$,
alors \ill{} \uncertain{peut} considérer le groupe \add{\uncertain{pro-fini}} analytique
\struck{$\pi$} $= G(\mathbb{Q}_\ell)$ ; \struck{et le groupe} \add{\ill{} $\pi$ \ill{}}
\ill{}, \uncertain{compact}. Alors \struck{\ill{} \ill{} \ill{} \ill{} considérer
le pro-groupe \ill{} \ill{} $\pi^{\mathrm{alg}}$ \ill{}, qui est}
\add{\uncertain{analytique}} \struck{\ill{}} \add{$\pi^{\mathrm{alg}}$} le groupe
\add{pro-algébrique (\uncertain{associé})} \ill{} une \emph{extension} de
$\pi$ \struck{\ill{}} (considéré comme un groupe pro-fini\add{$^{\mathrm{alg}}$}) par
\note{la phrase se poursuit p. 189}

\page{189}
\note{p. II.38 de l'auteur}
un \struck{\ill{}} groupe pro-algébrique connexe $(\widetilde{G}_{E'}^{0})_{E'}$,
\add{\uncertain{enveloppe} \uncertain{connexe} \ill{}} \struck{(\uncertain{enveloppe})}
\marginal{en haut à droite : \uncertain{Je n'ai} \ill{}}
\add{(\uncertain{algébrique})} algébriques \uncertain{des} représentations,
\struck{\ill{} \ill{}} $\ell$-adiques de $\pi$. Si p.ex. l'algèbre de Lie
$\mathfrak{g} = \pi^{\mathrm{an}} = \mathrm{Lie}(G)$\note{sic : « $\pi^{\mathrm{an}}$ » pour $\mathrm{Lie}^{\mathrm{an}}(\pi)$} est \uncertain{de} \uncertain{plus}
\uncertain{algèbre} \uncertain{dérivée}, \uncertain{dans} \uncertain{toutes} \uncertain{ses}
représentations $\ell$-adiques, l'image $G_{E'}$ de $\pi$ est \ill{} dans
$\widetilde{G}_{E'}$, donc $\widetilde{G}_{E'}^{0} \to \widetilde{G}_E^{0}$
\struck{\uncertain{induit}} \struck{\uncertain{un} \uncertain{isom.}} est un
isom. local au sens $\ell$-adique, \uncertain{donc} \struck{\uncertain{induit}} est
étale, \uncertain{donc} \struck{\uncertain{donc} \ill{} \uncertain{induit} \ill{}
\ill{} \ill{} $\widehat{G}^{0}$ \ill{} \ill{} \ill{} simplement connexe, \ill{}}
\add{\uncertain{l'unique} \uncertain{revêtement} \uncertain{universel}}
\add{\uncertain{simplement} connexe}, \ill{} \ill{} des \ill{}
\note{passage très raturé, encadré en partie et barré d'un trait sinueux}

\begin{tikzcd}[column sep=small, row sep=small]
\pi = G_{E'} \arrow[r, hook] \arrow[d] & \widetilde{G}_{E'} \arrow[d] \\
\pi = G_E \arrow[r, hook] & \widetilde{G}_E = G^{1}
\end{tikzcd}

\marginal{à droite du diagramme : s.-groupe alg. \ill{} de $G$ contenant $G^{0}$}

$G^{0}$ \struck{\uncertain{simple}} ½simple et simplement connexe, \ill{}
\ill{} \ill{} $G^{1}$ \uncertain{et} $\pi$ (\uncertain{alg.} : $G$ \uncertain{si}
\uncertain{que} l'enveloppe algébrique \ill{} $G$ est connexe) \uncertain{peut}
\uncertain{considérer} comme quotient d'une extension \uncertain{bien} déterminée de
$\pi$ (\uncertain{considéré} comme \ill{} groupe profini-algébrique) par
$G' = $ rev. \struck{universel} \add{simpl\supplied{ement} connexe} de $G^{0}$ ;
\ill{} \uncertain{donc} \ill{} un système projectif d'extensions de
groupes algébriques
\[
  1 \longrightarrow G' \longrightarrow G_\alpha \longrightarrow \pi_\alpha \longrightarrow 1
\]
\ill{} $\pi_\alpha$ \uncertain{parcourt} les quotients de $\pi$ par \uncertain{les}
s.-groupes \struck{\uncertain{distingués}} \add{\uncertain{ouverts}}
\add{\uncertain{suffisamment} \uncertain{petits}} \uncertain{\ill{} l'image de
$G' \to G^{0}$}, \add{(\uncertain{où} $G'/(G')^{0}$ \ill{})} \uncertain{évidemment}
$G^{1}$ \uncertain{étant} \uncertain{lui-même} \uncertain{un} \ill{} $G_\alpha$
\ill{} \uncertain{évidemment} \ill{} \add{\uncertain{des}} quotients $\pi_\alpha$ de
$\pi$ …).

On en conclut un élément canonique
\[
  \xi \in H^{2}(\pi, Z'(\mathbb{Q}_\ell))
\]
où $Z'$ est le \struck{\uncertain{noyau} \ill{}} \add{centre de} $G'$\struck{\ill{}},
l'annulation de $\xi$ étant nécessaire et suffisante pour que $\pi$ \uncertain{se}
\ill{} \ill{} un \uncertain{\ill{}} \uncertain{Borel} \ill{} \ill{}
\uncertain{connaissant} \uncertain{bien} \struck{\ill{}} \ill{} \ill{} \ill{}
\uncertain{vient} \uncertain{dans} $G'$ \add{en} \ill{} $H^{2}(\pi, Z'(\mathbb{Q}_\ell))$
\uncertain{dans} $H^{2}(\pi, Z(\mathbb{Q}_\ell))$ ($Z$ le centre \add{de $G^{0}$})
\struck{\uncertain{définit} \ill{}} est l'image de $\xi_0 \in H^{2}(G^{1}/G^{0}, Z(\mathbb{Q}_\ell))$,

\marginal{dans la marge gauche, verticalement : \ill{} \uncertain{défini} pour l'extension $G^{1}$ de $G^{1}/G^{0}$ par $G^{0}$}
\marginal{entre crochets, en marge : laissons tomber}

\page{190}
\note{p. II.39 de l'auteur}
6.9. \emph{Fonctorialité}. Un homom. de groupes profinis
\[
  u : \pi \longrightarrow \pi'
\]
définit
\begin{gather*}
  G = \pi^{\mathrm{an}} \to G' = \pi'^{\mathrm{an}}, \qquad
  \mathrm{Lie}(\pi^{\mathrm{an}}) = \mathfrak{g} \longrightarrow \mathrm{Lie}(\pi'^{\mathrm{an}}) = \mathfrak{g}' \\
  \widetilde{G} = \pi^{\mathrm{alg}} \to \widetilde{G}' = \pi'^{\mathrm{alg}}, \qquad
  \mathrm{Lie}(\pi^{\mathrm{alg}}) = \widetilde{\mathfrak{g}} \longrightarrow \mathrm{Lie}(\pi'^{\mathrm{alg}}) = \widetilde{\mathfrak{g}}'
  \end{gather*}
avec commutativité dans

\begin{tikzcd}[column sep=small, row sep=small]
G = \pi^{\mathrm{an}} \arrow[r] \arrow[d, "u^{\mathrm{an}}"] & \widetilde{G} = \pi^{\mathrm{alg}} \arrow[r] \arrow[d, "u^{\mathrm{alg}}"] & \pi \arrow[d, "u"] \\
G' = \pi'^{\mathrm{an}} \arrow[r] & \widetilde{G}' = \pi'^{\mathrm{alg}} \arrow[r] & \pi'
\end{tikzcd}

\begin{tikzcd}[column sep=small, row sep=small]
\mathfrak{g} = \mathrm{Lie}(\pi^{\mathrm{an}}) \arrow[r] \arrow[d] & \mathfrak{g}' = \mathrm{Lie}(\pi'^{\mathrm{an}}) \arrow[d] \\
\widetilde{\mathfrak{g}} = \mathrm{Lie}(\pi^{\mathrm{alg}}) \arrow[r] & \widetilde{\mathfrak{g}}' = \mathrm{Lie}(\pi'^{\mathrm{alg}})
\end{tikzcd}

Transitivités évidentes.

Si $u : \pi \to \pi'$ est un épimorphisme, il en est \uncertain{de} \uncertain{même}
de $u^{\mathrm{an}}$ et $u^{\mathrm{alg}}$ (dans un sens évident).
Si l'image de $u$ est de dimension finie\note{sic ; on attend « d'indice fini »}, alors
$u$ induit un épimorphisme $G^{0} \to G'^{0}$, et \uncertain{donc} des
épimorphismes $\mathfrak{g} \to \mathfrak{g}'$, $\widetilde{\mathfrak{g}} \to \widetilde{\mathfrak{g}}'$.

6.10. Soit $\Pi = (\pi_\alpha)$ un système projectif \add{filtrant} de groupes
profinis. En vertu de 6.9, il définit un système projectif filtrant de
systèmes projectifs filtrants de groupes analytiques, groupes algébriques,
alg. de Lie et alg. de Lie algébriques.
\begin{gather*}
  G = \Pi^{\mathrm{an}}, \quad \mathrm{Lie}(\Pi^{\mathrm{an}}) = \mathfrak{g} \\
  \widetilde{G} = \Pi^{\mathrm{alg}}, \quad \mathrm{Lie}(\Pi^{\mathrm{alg}}) = \widetilde{\mathfrak{g}} .
  \end{gather*}

\page{191}
\note{p. II.40 de l'auteur}
Ce dernier donne lieu aussi \uncertain{à} des homom.
\[
  \Pi^{\mathrm{an}} \longrightarrow \Pi^{\mathrm{alg}} \longrightarrow \Pi
\]
et
\[
  \Pi \simeq \Pi^{\mathrm{alg}}/(\Pi^{\mathrm{alg}})^{0}
\]
\[
  \mathfrak{g} = \mathrm{Lie}(\Pi^{\mathrm{an}}) \longrightarrow \widetilde{\mathfrak{g}} = \mathrm{Lie}(\Pi^{\mathrm{alg}}) .
\]
Fonctorialité pour $\Pi$ variable comme pro-objet de la catégorie des groupes
profinis.

\marginal{un grand trait en marge gauche}
Attention, $\Pi^{\mathrm{an}}$ et \struck{\ill{}} $\pi'^{\mathrm{an}}$ ne sont pas
pareils, si $\pi' = \varprojlim \pi_\lambda$. \struck{\uncertain{De} \uncertain{même}
pour $\Pi^{\mathrm{alg}}$, \ill{} \ill{} \ill{}} \struck{et les $\mathfrak{g}$,
$\widetilde{\mathfrak{g}}$. \ill{} \ill{} \ill{} épimorphismes de transition}

\emph{Exemple} 6.10.1. Les $\pi_\lambda$ finis, épimorphismes de transition
surjectifs, $\pi' = \varprojlim \pi_\lambda$ groupe analytique de dimension $> 0$.
Alors $\Pi^{\mathrm{an}}$ n'est autre que $(\pi_\lambda)$ lui-même, est formé de
groupes discrets, mais $\pi'^{\mathrm{an}}$ contient \add{« \uncertain{continûment} »}
$\pi'$ ! [\add{\uncertain{Ici}} $\Pi^{\mathrm{an}}$ \struck{plus} est « plus petit »
que $\pi'^{\mathrm{an}}$ !].

\emph{Exemple} 6.10.2. Les $\pi_\lambda$ les sous-groupes d'indice fini d'un $\pi$.
Alors $\pi' = \{e\}$, donc $\pi'^{\mathrm{an}}$ et $\pi'^{\mathrm{alg}}$ sont les
pro-groupes unités. Mais $\Pi^{\mathrm{an}}$ est le système des s.-groupes
\add{\uncertain{ouverts}} \struck{\uncertain{de} \uncertain{dimension} finie} des
$G_\alpha$ définissant $\pi^{\mathrm{an}}$, et $\Pi^{\mathrm{alg}}$ n'est autre que
$\widetilde{G}^{0}$, où $\widetilde{G} = \pi^{\mathrm{alg}}$. \struck{Donc} [Ici
$\Pi^{\mathrm{an}}$ est « plus grand » que $\pi'^{\mathrm{an}}$, qui est trivial !]

\page{192}
\note{p. II.41 de l'auteur}
6.11.
\[
  \mathrm{Mod}_{\mathbb{Q}_\ell}(\Pi^{\mathrm{an}}) \simeq \mathrm{Mod}(\Pi^{\mathrm{alg}}) \simeq \varinjlim_\lambda \mathrm{Mod}_{\mathbb{Q}_\ell}(\pi_\lambda)
\]
\note{le premier membre est vraisemblablement surchargé ; on attend $\mathrm{Mod}_{\mathbb{Q}_\ell}(\Pi)$}

6.12. Un cas important est celui où dans le \struck{système} pro-objet
$\Pi = (\pi_\lambda)$, les morphismes de transition sont quasi-surjectifs (i.e.
image d'indice fini). Alors il en est de même dans $\Pi^{\mathrm{an}}$ et
$\Pi^{\mathrm{alg}}$ (en particulier $(\Pi^{\mathrm{alg}})^{0}$ est un pro-objet
\emph{strict}). De plus, les syst. projectifs $\mathfrak{g} = \mathrm{Lie}(\Pi^{\mathrm{an}})$
et $\widetilde{\mathfrak{g}} = \mathrm{Lie}(\Pi^{\mathrm{alg}})$ sont aussi stricts ….

6.13. Soit
\[
  X = \varprojlim_\lambda X_\lambda
\]
un schéma, limite \add{($\mathfrak{X}$)} projective d'un système projectif filtrant de
schémas, à morphismes de transition affines. Soit $\xi$ un pt géom. de $X$,
$\xi_\lambda$ son image dans $X_\lambda$. Considérons
\[
  \Pi = \pi_1(\mathfrak{X}, \xi) = (\pi_1(X_\lambda, \xi))_\lambda ,
\]
système projectif de groupes profinis. Les représentations $\ell$-adiques de
\uncertain{ce} \uncertain{pro-objet} \struck{système} \uncertain{forment} \struck{\ill{}}
une catégorie équivalente à la catégorie
\[
  \varinjlim_\lambda (\text{Faisceaux \add{$\mathbb{Q}_\ell$-}\struck{\ill{}} tordus sur } X_\lambda),
\]
\uncertain{\ill{}} \uncertain{bien} \uncertain{entendu}. Ces catégories s'expliquent en
termes de $\Pi^{\mathrm{an}}$ ou $\Pi^{\mathrm{alg}}$ comme \uncertain{on} \uncertain{l'a}
\uncertain{vu} plus haut, et prennent ainsi (\uncertain{partiellement}) \uncertain{l'étude} :
l'aide de $\mathrm{Lie}(\Pi^{\mathrm{an}})$ et $\mathrm{Lie}(\Pi^{\mathrm{alg}})$.

\page{193}
\note{p. II.42 de l'auteur}
6.13.\note{numéro répété} Le cas le plus intéressant est celui où les $X_\lambda$ sont de
\emph{type fini} sur \struck{$X_\lambda$} $\mathbb{Z}$, ou sur un corps premier,
\struck{Alors} On sait que le foncteur \add{$\mathfrak{X} \to X$}
(\uncertain{limite}) des pro-objets de la catégorie des schémas de type fini sur
$\mathbb{Z}$, dans celle des schémas est pleinement fidèle, i.e. $\mathfrak{X}$ est
déterminé par la connaissance de $X$.

\emph{Ex} $X$ de présentation finie sur un schéma affine, p.ex. sur un corps.
En particulier $X$ \uncertain{local}, p.ex. spectre d'un corps, p.ex. d'un corps
alg. clos (\uncertain{sic})\note{lecture de la parenthèse douteuse}.

Dans ce cas, $\Pi$ mérite d'être \add{\uncertain{appelé}} \add{\uncertain{notation}
$\pi_1^{\mathrm{arith}}(X,\xi)$} le « \add{pro-}groupe fondamental
\emph{arithmétique} de $X$ », (on se rappellera que c'est un \struck{groupe}
pro-groupe topologique, \uncertain{admettant} un \uncertain{homom.} canonique
\[
  \pi_1(X,\xi) \xrightarrow{\ \sim\ } \varprojlim \Pi = \varprojlim_\lambda \pi_1(X_\lambda, \xi) ,
\]
mais on se rappellera que la connaissance de $\pi_1(X,\xi)$ ne détermine
\uncertain{pas} \uncertain{celle} du \add{\uncertain{pro-objet}}
\struck{\ill{}} $\Pi = \Pi_1(X,\xi)$. Par exemple, si $X$ est le spectre d'un
corps alg. clos, $\pi_1(X,\xi) = e$, tandis que $\Pi$ n'est \uncertain{alors}
jamais trivial, \uncertain{\ill{}} \uncertain{nous} \uncertain{allons} \uncertain{voir}.

En vertu de 6.10, $\pi^{\mathrm{arith}}(X,\xi)$ \ill{} \ill{}

\page{194}
\note{p. II.43 de l'auteur}
\uncertain{naissance} à \struck{des} \uncertain{des} groupes pro-analytiques et
pro-algébriques
\[
  \pi_1^{\mathrm{arith}}(X,\xi)^{\mathrm{an}}, \quad \pi_1^{\mathrm{arith}}(X,\xi)^{\mathrm{alg}}
\]
et à des pro-algèbres de Lie correspondantes.

6.14. Supposons $X$ intègre et normal, alors on peut supposer les $X_\lambda$
intègres, normaux (quitte à \struck{\ill{}} remplacer $X_\lambda$ par le normalisé
de l'image schématique \uncertain{fermée} de $X$ dans $X_\lambda$ …).
\struck{\ill{}} Comme tout morphisme dominant de t.f. de schémas normaux,
$X_\mu \to X_\lambda$ induit un \struck{\ill{}} homom. \emph{quasi-surjectif}
\[
  \pi_1(X_\mu, \xi_\mu) \longrightarrow \pi_1(X_\lambda, \xi_\lambda),
\]
i.e. on est dans le cas de 6.12. En particulier, les \struck{\ill{}} homom. de
transition des \struck{\ill{}} $G = \pi_1^{\mathrm{arith}}(X,\xi)^{\mathrm{an}}$ et
$\widetilde{G} = \pi_1^{\mathrm{arith}}(X,\xi)^{\mathrm{alg}}$ sont des
quasi-épimorphismes, et les homom. de transition \struck{\ill{}} \uncertain{dans}
$\mathfrak{g}$ et $\widetilde{\mathfrak{g}}$ sont des épimorphismes (\uncertain{donc}
\uncertain{les} $\mathfrak{g}$, $\widetilde{\mathfrak{g}}$ sont des pro-objets
\emph{stricts}).
\marginal{encadré, rattaché au paragraphe suivant : si $\ell$ est inversible sur $X$ (\uncertain{sauf} \uncertain{mention} \uncertain{contraire} …) \uncertain{alors}}

6.15. Notons que \struck{si $\ell$ est inversible sur $X$} \add{\ill{}}
\struck{\ill{}} \add{\uncertain{donc}} $\Pi = \pi_1^{\mathrm{arith}}(X,\xi)$ n'est
\struck{\ill{}} \add{\uncertain{jamais}} « discret » i.e. \ill{} $\dim \mathfrak{g} > 0$,
\struck{\ill{} $\mathfrak{g}$} \uncertain{ou} \uncertain{encore} $\dim \widetilde{\mathfrak{g}} > 0$.
(En fait, on montre que en \uncertain{\ill{}} dimensions sont \emph{très} \uncertain{grandes}).

En effet, il suffit de \struck{\ill{} \ill{} \ill{} \ill{}} \add{\uncertain{voir} \uncertain{que} \ill{}}
\uncertain{que} si $X$ est de type fini sur $\mathrm{Spec}\,\mathbb{Z}$, \ill{} \ill{}
d'un \add{\uncertain{tel} \uncertain{que} $\ell$ \uncertain{y} soit inversible,}
\uncertain{point} \uncertain{géom.} $\xi$, \uncertain{alors}, \struck{$\pi_1(X$,}
\uncertain{posant} $\pi = \pi_1(X,\xi)$, $\pi^{\mathrm{alg}}$ \uncertain{n'est}
\uncertain{jamais} discret. Or

\page{195}
\note{p. II.44 de l'auteur}
\struck{Si $X$ est de car. $p$, alors $X$ est \ill{}.}
Si $X$ est de car. $p > 0$ ($p \neq \ell$), alors $X \to \mathrm{Spec}\,\mathbb{F}_p = S$
induit \struck{\ill{}} quasi-isom. $\pi_1(X,\xi) \to \pi_1(S,\xi) \simeq \hat{\mathbb{Z}}$,
donc \struck{\ill{}} $\pi_1(X,\xi)^{\mathrm{alg}} \to (\hat{\mathbb{Z}}^{\mathrm{alg}})$ est un
quasi-\add{épim.}\struck{\ill{}}, et $\hat{\mathbb{Z}}^{\mathrm{alg}}$ (\uncertain{comme}
\uncertain{on} \uncertain{voit} \uncertain{aussitôt}) est de dim. infinie. Si $X$ est de
car. $p = 0$, alors \uncertain{posons} $S = \mathrm{Spec}\,\mathbb{Z}[1/\ell]$,
\struck{\ill{} $\pi = \pi_1(X,\xi) \to \pi_1(S,$} \uncertain{le} raisonnement
\uncertain{est} le même, en notant que $\pi_1(S)$ a \struck{un quotient} comme
\uncertain{plus} grand quotient abélien $\mathbb{Z}_\ell^{*}$, qui est
quasi-isomorphe à $\mathbb{Z}_\ell$, or $\mathbb{Z}_\ell^{\mathrm{alg}}$ est encore de
dim. infinie ….

6.16. Revenons aux notations générales de 6.13. L'intérêt de la considération
\struck{\ill{}} \uncertain{de} $\Pi = \pi_1^{\mathrm{arith}}(X,\xi)$ tient au fait que
tous les faisceaux $\ell$-adiques \uncertain{constructibles} \struck{\ill{}}
\uncertain{tordus} qui \uncertain{interviennent} pratiquement sont de
\uncertain{nature} \uncertain{arithmétique}, i.e. proviennent déjà d'un tel faisceau
sur un $X_\lambda$, \uncertain{est} unique modulo \uncertain{passage} à un $X_\mu$, avec
$\mu \gg \lambda$. \struck{En \ill{}} \struck{\ill{},} \ill{} \uncertain{de} \uncertain{même}
\uncertain{pour} les \uncertain{homom.} de faisceaux $\ell$-adiques. En d'autres termes,

\page{196}
\note{p. II.45 de l'auteur}
il y a lieu de travailler dans la catégorie
\[
  \varinjlim_\lambda (\text{faisceaux tordus sur } X_\lambda),
\]
qui précisément s'explicite en termes de représentations \add{$\ell$-adiques}
\struck{\uncertain{linéaires}} du pro-groupe $\Pi$, \uncertain{ou} en termes de
$\Pi^{\mathrm{alg}}$, i.e. de repr. alg. de $\Pi^{\mathrm{alg}}$.

\emph{Exemple 1}. Soit $f : Y \to X$ un morphisme propre et lisse, on peut
\uncertain{prendre} les $R^{i}f_*(\mathbb{Q}_{\ell,Y})$, \uncertain{\ill{}}
\struck{\ill{} \uncertain{donne}} les $R^{i}f_*(\mathbb{Q}_{\ell,Y})$ \uncertain{définissent}
une représentation de $\widetilde{G}$, \uncertain{ou} $\widetilde{\mathfrak{g}}$. De même,
les homom. entre $R^{i}f_*$ définis par des $X$-morphismes $Y \to Y'$, ou
[\uncertain{correspondances} \struck{\ill{}} \uncertain{définies} \uncertain{sur} $X$
et intègres \uncertain{relat.}) par des \struck{\ill{}} \struck{\ill{}} \uncertain{cycles}
\add{algébriques} \uncertain{correspondants} \add{\uncertain{sur}} \uncertain{les} fibres
génériques.

\emph{Exemple 2} (\uncertain{Cent} \uncertain{élémentaire}) : $T_\ell(A)$, où $A$
est un schéma abélien sur $X$, \struck{$= T_\ell(A) \otimes_{\mathbb{Z}_\ell} \mathbb{Q}_\ell$}
\add{(\uncertain{Noter} que}
\[
  \underline{\mathrm{Hom}}(T_\ell(A), \mathbb{Z}_\ell) \simeq R^{1}f_*(\mathbb{Z}_{\ell,A}),
\]
donc $T_\ell(A)$ et $R^{1}f_*(\mathbb{Z}_{\ell,A})$ se déterminent \uncertain{mutuellement}.
\struck{\ill{}}

\uncertain{Pour} \uncertain{que} \struck{\ill{}} \uncertain{cela} \uncertain{définisse}
une représentation de $\mathfrak{g} = \mathrm{Lie}(\pi_1^{\mathrm{arith}}(X;\xi))$, il faut
\uncertain{seulement} \uncertain{passer} \uncertain{à} \uncertain{l'extension} par
$\mathbb{Q}_\ell$.\note{fin de page très incertaine}

\page{197}
\note{p. II.46 de l'auteur}
Bien entendu, les exemples 1, 2 \add{\uncertain{sont} \uncertain{des} \uncertain{repr.} de
$\pi^{\mathrm{an}} \to \pi^{\mathrm{alg}}$} \uncertain{dépendent} \uncertain{de}
\uncertain{façon} \emph{fonctorielle} \uncertain{pour} $X$, \uncertain{resp.} $A$,
variables. En particulier, dans le cas d'un $A$, on trouve un homom. d'anneaux
(injectif)
\[
  (6.15.1)\qquad \mathrm{End}(A) \otimes_{\mathbb{Z}} \mathbb{Q}_\ell \hookrightarrow
  \mathrm{End}_{\Pi}\bigl(\underbrace{T_\ell(A(\xi)) \otimes_{\mathbb{Z}_\ell} \mathbb{Q}_\ell}_{E}\bigr) .
\]

\emph{Remarque 6.16}. Une des conséquences des \emph{conjectures de Tate}
(qui montrent le rôle essentiel des représentations \uncertain{\ill{}} du groupe
fondamental \struck{\ill{}} arithmétique) est que \uncertain{le} homom. précédent
est bijectif \struck{\ill{}} lorsque $X$ est normal intègre. Pour cette question,
on se ramène d'ailleurs aussitôt au cas où $X$ est le spectre d'un corps
\emph{algébriquement clos}. \uncertain{Cela} montre à quel point (\uncertain{même}
dans un cas \uncertain{réputé} « géométrique ») l'opération de $\Pi$ est loin d'être
triviale. [C'est prouvé par Tate si $X$ est le spectre d'un corps fini,
\struck{\ill{} de façon équivalente, si \ill{} \ill{}} \uncertain{et} \uncertain{si}
\uncertain{c'est} \uncertain{le} spectre de la clôture algébrique d'un corps fini ….].

\emph{Remarque 6.17}. \uncertain{On} \uncertain{s'est} \uncertain{rendu} compte que
$\Pi$ lui-même, ou $G = \Pi^{\mathrm{an}}$, $\widetilde{G} = \Pi^{\mathrm{alg}}$,
$\mathfrak{g}$, $\widetilde{\mathfrak{g}}$, sont \uncertain{démesurément} grands et ils
échappent pratiquement à notre contrôle. Leur rôle est surtout théorique, via
leurs représentations $\ell$-adiques

\page{198}
\note{p. II.47 de l'auteur}
Dans chaque \uncertain{problème} géométrique particulier, on n'a guère qu'à
examiner les représentations (\uncertain{analytiques} ou algébriques) se
\struck{\uncertain{factorisant}} \uncertain{par} un groupe \struck{quotient}
$G_\alpha$ resp. $\widetilde{G}_\alpha$ \uncertain{fixé} du syst. projectif, dont on
peut supposer que c'est l'image dans un $\underline{\mathrm{Aut}}_{\mathbb{Q}_\ell}(E)$
d'un groupe $\pi_1(X_\lambda, \xi_\lambda)$, resp. son enveloppe algébrique.
Notons que \struck{\ill{}} dans (6.15.1) l'action de $\Pi$ sur $E$ définit un
groupe algébrique $\widetilde{G}_E \subset \underline{\mathrm{Aut}}_{\mathbb{Q}_\ell}(E)$,
enveloppe algébrique d'un groupe analytique $G_E$ \struck{\ill{}} \add{image} d'un
$\pi_1(X_\lambda, \xi)$, et que le deuxième membre de (6.15.1), \uncertain{peut}
aussi \uncertain{s'écrire}, \uncertain{tautologiquement},
\[
  \mathrm{End}_{G_E}(E) = \mathrm{End}_{\widetilde{G}_E}(E).
\]
Si $\pi_1(X,\xi)$ (\struck{\ill{}} \uncertain{arith.}) $= 1$ (p.ex. $X$ le spectre
d'un corps alg. clos) alors le deuxième membre \uncertain{s'écrit} aussi
\[
  \mathrm{End}_{\mathfrak{g}_E}(E) = \mathrm{End}_{\widetilde{\mathfrak{g}}_E}(E),
\]
où $\mathfrak{g}_E = \mathrm{Lie}(G_E) \subset \mathrm{End}(E)$ et $\widetilde{\mathfrak{g}}_E$
est l'enveloppe algébrique de $\mathfrak{g}_E$ ….

\note{la rédaction s'arrête ici, en milieu de page ; le chapitre II ne se poursuit pas dans ce lot}

\section{SGA 7 \uncertain{IX} appendice. Applications des conjectures standard sur les cycles algébriques}
\note{titre écrit de sa main sur une feuille de couverture (p. 199 du fonds), sans autre contenu ; le numéro romain est douteux}

\page{200}
1) \emph{Cohomologie des diviseurs à croisements normaux.}
\[
  Y = \bigcup_\alpha Y_\alpha \qquad Y_\alpha \text{ projectif et lisse, } Y_{\alpha_0} \cap \dots \cap Y_{\alpha_p} \text{ est lisse.}
\]
\note{sous $Y$, un trait et le mot « propre »}
\[
  H^{n}(Y, \mathbb{Q}_\ell) \Longleftarrow E_2^{pq} = H^{p}(\sigma \longmapsto H^{q}(Y_\sigma, \mathbb{Q}_\ell)), \qquad
  Y_{\alpha_0 \dots \alpha_p} = Y_{\alpha_0} \cap \dots \cap Y_{\alpha_p}
\]
\struck{\ill{}} $E_2^{pq}$ est de \emph{poids} $q$, donc $d_r^{pq}$ est nul pour \uncertain{tout} $r$.
$H^{i}(Y, \mathbb{Q}_\ell)$ se divise en les $E_2^{p,q}$ ($p+q = i$), de poids $0, 1, \dots, i$.

\marginal{en biais, à l'encre brune : NB des plus \uncertain{besoin} (conjectures standard) \ill{} les dimensions des $E_2^{pq}$ \uncertain{indépendantes} de $\ell$}

2) \emph{Cohomologie des \struck{variétés} \add{schémas} quasi-projectifs lisses non propres.}
\[
  X = X' - Y, \qquad X' \text{ projectif et lisse, } Y \text{ diviseur}
\]
à croisements normaux, tels que les comp. irréd. $Y_i$ soient lisses.
\begin{gather*}
  \text{\struck{$H^{i-1}(X', \mathbb{Q}_\ell)$}} \to H^{i-1}(Y, \mathbb{Q}_\ell) \to H^{i}_{!}(X, \mathbb{Q}_\ell) \to H^{i}(X', \mathbb{Q}_\ell) \\
  \to H^{i}(Y, \mathbb{Q}_\ell) \to H^{i+1}_{!}(X, \mathbb{Q}_\ell) \to \cdots
  \end{gather*}
\note{marqué (a) dans un cercle en marge ; sous $H^{i-1}(Y)$ « poids $\leqslant i-1$ », sous $H^{i}(X')$ « poids $i$ », sous $H^{i}(Y)$ « poids $\leqslant i$ »}

On trouve que
\[
  \text{a) \struck{l'image}}\ \mathrm{Im}\bigl(H^{i-1}(Y, \mathbb{Q}_\ell) \to H^{i}_{!}(X, \mathbb{Q}_\ell)\bigr)
  = \mathrm{Ker}\bigl(H^{i}_{!}(X, \mathbb{Q}_\ell) \to H^{i}(X', \mathbb{Q}_\ell)\bigr)
\]
\emph{ne dépend pas} de la résolution choisie : c'est la \uncertain{partie}
\uncertain{de} poids $\leqslant i-1$ de $H^{i}_{!}(X, \mathbb{Q}_\ell)$.

\struck{\ill{}} Les Ker et Coker de
\[
  \varphi^{i} : H^{i}(X', \mathbb{Q}_\ell) \to H^{i}(Y, \mathbb{Q}_\ell),
\]
à \uncertain{isom.} \uncertain{cano.} près, \uncertain{sont} aussi indépendants de la
résolution choisie. \struck{\ill{}} \add{\uncertain{Considérons} la filtration de
\uncertain{division}} \struck{\ill{}} \ill{} $H^{i}(Y, \mathbb{Q}_\ell)$ provenant de la
suite spectrale étudiée dans 1),
\[
  E_2^{pq} = H^{p}(\sigma \mapsto H^{q}(Y_\sigma, \mathbb{Q}_\ell)),
\]
correspondant aux \uncertain{\ill{}} \uncertain{des} $(p+q = i)$. L'homom.
\[
  H^{i}(Y, \mathbb{Q}_\ell) \to E_2^{0,i} = H^{0}(\sigma \mapsto H^{i}(Y_\sigma, \mathbb{Q}_\ell)) = \mathrm{Ker}\bigl(H^{i}(Y_\alpha, \mathbb{Q}_\ell) \rightrightarrows \cdots\bigr)
\]
\note{le dernier terme, surchargé en bas de page, est lu d'après le contexte ; la phrase se poursuit au-delà de ce lot}

\end{document}
